\chapter*{Önsöz}
\addcontentsline{toc}{chapter}{Önsöz}
\markboth{Önsöz}{Önsöz}
%  SAYFA DOLULUĞU (Önsöz): Bu metin 2. sayfayı SON SATIRINA KADAR dolduruyor.
% Teşekkür paragrafına metin EKLERSENİZ son satır 3. sayfaya düşer (yetim satır).
% Kural: eklenen kadar kırpın (net satır sayısı sabit kalsın) veya \enlargethispage'i
% artırın — ama en fazla 1 satır; 2 satır sayfa numarasını kenara yapıştırır.
% Denetim: python3 scripts/check_frontmatter.py
\enlargethispage{\baselineskip}

Lisedeki son iki yılımda alt dönemlerimden en çok duyduğum cümle, \textit{``Abi sınava az kaldı ve ne yapmam gerektiğini bilmiyorum''} oldu. Bu soruyu soranlar tembel ya da isteksiz değil, aksine son derece parlak öğrencilerdi. 1. Aşama Sınavı kendine has yapısıyla özel bir hazırlık gerektirirken öğrencilerin önünde ne doğru dürüst bir yol haritası ne de doğrudan bu sınava odaklanan derli toplu bir kaynak vardı. Yıllarca bu arkadaşlarıma en fazla bir-iki ay yetecek dağınık notlar önerebilmenin burukluğunu yaşadım.

Bilgisayar olimpiyatı salt kod yazmaktan ibaret değildir; bir problemi matematiksel olarak modelleme ve bulunan fikri hayata geçirme işidir. Başarılı olanların ortak noktası soru kalıplarını ezberlemek değil, çözdükleri her sorudan bir ders çıkarıp bunu bir sonraki probleme uyarlayabilmeleridir. Bu kitap da yarışma hilelerinden ziyade sonlu matematik, sayılar teorisi ve algoritmaların kavramsal temellerini sağlam atmayı hedefler.

Bilgisayar olimpiyatı ağır bir zihinsel spordur. Saatlerce tek bir soruyu çözemediğiniz, bazı soyut konuları günlerce kafanızda oturtamadığınız zamanlar olacaktır. Önemli olan bu zorlanma anlarında pes etmemek ve kavramların zihninizde yer etmesi için kendinize zaman tanımaktır. Doğru bir yönlendirme ve düzenli çalışmayla geliştirilen problem çözme sezgisi, ileride hangi alana yönelirseniz yönelin karşılığını fazlasıyla verecektir.

Üniversiteye başlamadan önceki yaz tatilimde, henüz lisedeki tecrübelerim tazeyken bu kitabı hazırlamak istedim. Günümüzdeki yapay zekâ araçları dizgi ve kaynak tarama gibi rutin işleri oldukça kolaylaştırdı; ben de zamanımı müfredatı kurgulamaya, soruları titizlikle ayıklamaya ve adımları kontrol etmeye ayırabildim. Yaklaşık 600 sayfa olan bu kitabı acele etmeden, örnekleri kendi kâğıdınızda çözerek ve zamana yayarak sindirmenizi öneririm.

Herkesten çok destek olan aileme; birlikte saatlerce soru çözdüğümüz okul arkadaşlarıma, alt dönemlerime, TÜBİTAK Fen Lisesi öğretmenlerime ve okul müdürüm Dr. Hacı Ahmet Yıldırım'a teşekkür ederim. Destekleri için TÜBİTAK ve BİDEB'e, TÜBİTAK Bilgisayar Olimpiyatları komite başkanı Prof. Dr. Muhammed Fatih Demirci'ye, kamplardaki rehberlerime, desteğini hissettiğim Ahmet Kaan Avcı ağabeyime, özverili BİDEB personeli Dr. Burcu Çalışkan'a, üretkenlik konusunda bana rol model olan Dr. Fatih Kürşat Cansu'ya, beni bu projeye başlamaya YouTube'daki videolarıyla cesaretlendiren Prof. Dr. Oğuz Ergin'e ve derslerime katılan öğrencilerime şükranlarımı sunarım.

\clearpage